\documentclass{situationist-preprint}

\title{The B-Series Suffices}
\msid{TSM-01-c}
\msfamily{TSM-01}
\msdiscipline{Time, Space and Modality}
\mstype{Consequence}
\mskeywords{B-theory; tense; change; direction of time; open future; self-locating belief}

\begin{document}
\maketitle

\begin{abstract}
TSM-01 concludes that the A-series is incoherent and that time is real and constituted by the B-series, the causal and metrical structure of spacetime. A common worry is that a B-series alone leaves out something time must supply: change, direction, the truth of tensed sentences, the openness of the future, or the significance of the present. We derive from TSM-01, with results of neighboring families, that each of these is supplied by the B-series together with the position of thinkers in it. Nothing temporal requires an addition to the B-series, and so the incoherence of the A-series costs nothing.
\end{abstract}

\section{What time is asked to do}

Five roles are standardly assigned to time: (i) to order events and measure their separation; (ii) to make change possible; (iii) to give the order a direction; (iv) to make tensed sentences and thoughts true or false; and (v) to leave the future open while the past is fixed. A sixth is assigned by experience: (vi) to present itself as passing. An A-theorist holds that at least one of these requires a privileged and moving present. Given TSM-01, that would leave the role unfilled.

\begin{definition}[Sufficiency]\label{d:suff}
The B-series \emph{suffices} for a role if the facts that fill it are fixed by the B-facts (what occurs at each time and the relations of earlier, later and simultaneous) together with facts about where thinkers and their thoughts are located among them.
\end{definition}

The second clause adds no tensed facts. That a thought occurs at a given time is itself a B-fact.

\section{The derivation}

\begin{proposition}\label{p:main}
The B-series suffices for roles (i)--(vi).
\end{proposition}

\begin{proof}
(i) Order and duration are the causal and metrical structure of spacetime, which is what TSM-01 identifies the B-series with. (ii) Change is an object's having incompatible properties at different times \citep[§442]{Russell1903}; this is a B-fact, and TSM-01-b argues that the demand for any other kind of change rests on an equivocation. (iii) The direction of time is the orientation of the physical asymmetries along the B-series, which TSM-05 derives from the low entropy of the initial state; it is fixed by what occurs at each time. (iv) A tensed token is true if and only if it stands in the appropriate B-relation to the event it concerns: \enquote{it is raining now}, uttered at $t$, is true just if it rains at $t$ \citep{Mellor1998}. That a time is present is that a thought or utterance occurs at it (TSM-15), and \enquote{now} is an indexical (TSM-18). (v) The B-series may branch: relative to each moment there is one past and several possible continuations, and a future contingent is settled at a moment if it holds on every continuation through it \citep{Belnap2001}. Settledness is a relation between a moment and a sentence, so openness requires no privileged moment (TSM-12). (vi) The impression of passage is fixed by the B-facts about experience (TSM-02).
\end{proof}

\begin{corollary}\label{c:nocost}
Given TSM-01, the incoherence of the A-series removes no temporal role.
\end{corollary}

Two consequences deserve stating. The first concerns ontology. Since the B-series is the same structure at every moment, no moment has an ontological status that others lack; TSM-03 accordingly holds that the past of each moment exists, and the growing block, which privileges an edge, is not available to a view on which the B-series suffices. Branching does not reintroduce an edge: in a branching B-series each moment has many continuations, and none of them is the one that \enquote{will come to exist}, since there is no moving present for it to come to exist at. The second concerns knowledge. What we know in knowing what time it is, is self-locating information, our position in a B-series that we otherwise know; EPI-18-c and REL-19-b develop this for credence and for divine omniscience.

\section{Objections and replies}

\begin{objection}[Relief]
\citet{Prior1959} observes that one may say \enquote{thank goodness that's over} of a headache, and that what one is thankful for is not that the headache ends before this utterance, a B-fact one could be indifferent to.
\end{objection}

\begin{reply}
What one is thankful for is self-locating: that one's present thinking is later than the headache. Its content is true at some centres and false at others, as EPI-18-c argues, and it is the content one has reason to welcome, since it is the location of one's present thinking that determines whether one is in pain. Role (iv) requires nothing more \citep{Mellor1998}.
\end{reply}

\begin{objection}[Becoming]
In a branching structure the future becomes settled, and becoming is A-change under another name.
\end{objection}

\begin{reply}
A sentence is settled at one moment and unsettled at an earlier one. That is a variation across moments, a B-fact about the branching structure, and not a movement of settledness along it. To say that a contingent \enquote{has now become} settled is to say that it is settled at the moment of utterance (TSM-12, TSM-18).
\end{reply}

\section{Conclusion}

Order, change, direction, tensed truth, the open future and the experience of passage are all fixed by the B-series and the location of thinkers in it. The B-series therefore suffices, and the tenseless view of TSM-01 gives up none of what time was asked to do.

\begin{thebibliography}{9}
\bibitem[Belnap et~al.(2001)]{Belnap2001} Belnap, N., Perloff, M., and Xu, M. (2001). \emph{Facing the Future: Agents and Choices in Our Indeterminist World}. Oxford University Press.
\bibitem[Mellor(1998)]{Mellor1998} Mellor, D.~H. (1998). \emph{Real Time II}. Routledge.
\bibitem[Prior(1959)]{Prior1959} Prior, A.~N. (1959). Thank goodness that's over. \emph{Philosophy}, 34(128), 12--17.
\bibitem[Russell(1903)]{Russell1903} Russell, B. (1903). \emph{The Principles of Mathematics}. Cambridge University Press.
\end{thebibliography}

\end{document}
